\section{Introduction}
\label{sec:intro}

The public ALMA archive holds thousands of hours of millimetre and
submillimetre observations taken for other reasons, mostly studies of dust
discs and planet formation. Their visibilities allow a spectrum at the
position of the star at the centre of each field, at a resolution and depth
set by the original observer's correlator setup. This paper searches those
spectra for a narrowband transmitter at the star's own position. Nothing of
the kind survived.

Searches for \emph{technosignatures} --- detectable imprints of technology on
the electromagnetic spectrum, intended or not
\citep{Sheikh2020,SETIreview2025} --- have examined many thousands of stars
over six decades, yet the last published accounting of that effort covers
only \HaystackFrac{} of the space of transmitter position, power, frequency,
bandwidth, polarisation and repetition rate
\citep{Wright2018}, and essentially all of it lies below
${\sim}30$\,GHz. No published search above that frequency has reached a star
within 40\,pc, though more distant stars have been searched in the
millimetre band \citep{Mason2025}. The band is attractive because a carrier
arrives less distorted and a given aperture delivers more gain per watt
supplied; what makes it hard belongs to the archive rather than to the sky,
and \S\ref{sec:background} sets that out.

The signal sought is an \emph{unresolved spectral carrier}: emission confined
to a few channels, unresolved at the stellar position, and present at more
than one epoch. Each qualifier is a restriction, and the last is the
strictest: it is recurrence rather than brightness that separates an
interesting event from a detection \citep{Sheikh2021}. The archive holds a second,
separated epoch for some of these systems and not for others, so what could
be searched and what could be confirmed are different extents. The frequency is
allowed to drift, as it must for a transmitter carried on a rotating,
orbiting body (\S\ref{sec:statistic}). ``Narrowband'' here describes the
assumed transmitter and never the resolving power of the archived
channels. The targets
were chosen by other programmes for other purposes, so every statistical
statement that follows is about an archive-selected sample
(\S\ref{sec:excludes}).
